In this section, I present some applications that can immediately benefit from this framework outside of software systems and also some future thoughts for interesting applications that I would want to explore within the world of software systems.

\subsection{Universal Design Language and Digital Twins}

Designs are the universal structures used across many disciplines. The ability to unambiguously establish a closed semantic world and reason about it has applications well beyond software systems. If implemented with the right stratergy, this framework has the ability to become a universal design language used in many disciplines.

Since the infromation contained in the DAL has unambiguous meaning, it can be used to establish a powerful visualization framework. This would allow engineers to visualize the closed semantic world in a custom way and gain insights into their design by reasoning about its meaning. More importantly, since this framework extends the design into a computable semantic model (CSM), it extends the CSM to become a digital twin of what it is modelling.

In software systems, the design establishes a closed semantic world that mirrors what partcipates in reality because the design is what defines it (this statement is even more valid with CSM's construted using CSP's). This is a unique property of software systems since they operate entirely on a stack that is symbollically defined by humans. However, in other disciplines, reality is modelled by engineers to be able to represent it accurately and predict it. In these applications, the CSM becomes the framework through which that modelling happens. 

In the digital twin models, information collected by sensors is used to inform the model and refine it so that it is a more accurate representation of reality and to improve its ability to make more accurate predictions. In these cases, the domain structure of the collected data is unambiguously established by the CSM, as such, the data can be domain compressed using CLP's engine. This extends the use of CSM and CLP beyond software systems. 

The applications of CSM are only limited by what can be modelled, anything that can be modelled as a closed semantic world can be reasoned about using the engine. A personal favourite of mine is to model concepts using the CSM so that the model can be used as a teaching tool. This applications exploits the fact that when humans teach, they are ultimately building a closed semantic world to convey the concept to another human mind. Using the CSM, the semantic world necessary to represent the idea can be transformed into a computable model in which every step in the process is informed and validated by the model. There are many interesting ways in which this can be used to turn teaching and learning into an unambiguous process that converges on complete understanding.

\subsection{Future Exploration: Automated Development}

The thoughts in this section are more future oriented but something that I would enjoy exploring. I believe that the Computable Semantic Model built from Computable Semantic Primitives opens the door to interesting approaches for automating the development of software systems. This is because the engine has the ability to reason about the world that was developed to identify its internal consistency and correctness. I think there are interesting techniques where a semantic toolkit is used to construct closed semantic worlds and expand the meaning of world.

Since failure diagnosis is automated and the design can deterministically learn new semantics through root cause analysis, I think there are applications where the design operates entirely in a virtual environment and integrates its meaning into the existing closed semantic world. In a solution like this the simulated environment itself has unambiguous meaning and when the design fails, it learns how it must modify its behavior to achieve its intentions in its environment.

It would be able to perform root cause analysis on the failure because the entire world is modelled and the meaning is unambiguous. While there are certainly many intermediate steps on the way to reaching this, I feel that when software systems are modelled in this way to establish its meaning as a computable model and they start to work together in a simulated world, solutions like this will become more realistic where designs try and fail in simulations to learn. 